\documentclass[12pt,a4paper]{article}

% ------------------------------------------------------------
% Packages
% ------------------------------------------------------------
\usepackage[utf8]{inputenc}
\usepackage[T1,T5]{fontenc}
\usepackage[english]{babel}
\usepackage[a4paper,margin=2.5cm]{geometry}
\usepackage{amsmath,amssymb}
\usepackage{graphicx}
\usepackage{booktabs}
\usepackage{array}
\usepackage{float}
\usepackage{enumitem}
\usepackage{listings}
\usepackage{xcolor}
\usepackage{fancyhdr}
\usepackage{hyperref}
\usepackage{microtype}

% ------------------------------------------------------------
% Editable report information
% ------------------------------------------------------------
\newcommand{\coursename}{Digital Signal Processing}
\newcommand{\coursecode}{CO2036}
\newcommand{\semester}{HK261}
\newcommand{\labtitle}{LAB 4 -- Z-TRANSFORM AND CONVOLUTION}
\newcommand{\studentone}{Nguyễn Văn Hoàng Phát -- 2352896}

% ------------------------------------------------------------
% Formatting
% ------------------------------------------------------------
\hypersetup{
    colorlinks=true,
    linkcolor=black,
    urlcolor=blue,
    citecolor=black
}

\pagestyle{fancy}
\fancyhf{}
\lhead{\coursename}
\rhead{Lab 4}
\cfoot{\thepage}
\setlength{\headheight}{15pt}

\setcounter{secnumdepth}{3}
\setcounter{tocdepth}{2}

% ------------------------------------------------------------
% Scilab listing style
% ------------------------------------------------------------
\lstdefinelanguage{Scilab}{
  morekeywords={
    function,endfunction,if,then,else,end,for,while,break,continue,
    select,case,otherwise,return,exec,clear,clc,disp,mprintf,
    zeros,ones,length,min,max,size,sum,abs,floor,fft,scf,clf,subplot,
    plot,plot2d3,xgrid,xtitle,error,wavread,playsnd,atomsLoad,imread,
    imshow,rgb2gray,histplot,matrix,double,pmodulo
  },
  sensitive=true,
  morecomment=[l]{//},
  morestring=[b]"
}

\lstset{
    language=Scilab,
    basicstyle=\ttfamily\small,
    keywordstyle=\bfseries,
    commentstyle=\itshape,
    stringstyle=\ttfamily,
    numbers=left,
    numberstyle=\tiny,
    stepnumber=1,
    numbersep=8pt,
    frame=single,
    breaklines=true,
    showstringspaces=false,
    tabsize=4,
    captionpos=b
}

\newcommand{\reportfigure}[4][0.92\textwidth]{%
\begin{figure}[H]
    \centering
    \includegraphics[width=#1]{#2}
    \caption{#3}
    \label{#4}
\end{figure}
}

\begin{document}

% ============================================================
% Cover
% ============================================================
\begin{center}
    {\large \textbf{\coursename\ -- \semester\ -- \coursecode}}\\[1.2cm]

    {\LARGE \textbf{DIGITAL SIGNAL PROCESSING LAB}}\\[1cm]
    {\Large \textbf{\labtitle}}\\[1.5cm]

    \begin{tabular}{>{\bfseries}l l}
        Student: & \studentone \\
    \end{tabular}
\end{center}

\vspace{1.2cm}

\section*{Overview}
This report studies the Z-transform and region of convergence (ROC) of
several discrete-time signals and investigates practical signal-processing
operations in Scilab. The experimental part includes audio analysis,
image processing using the IPCV toolbox, linear convolution using both
folding-and-shifting and matrix methods, and circular convolution using
the same two approaches. Numerical results, graphical outputs, and
energy calculations are used to verify the implementations.

\tableofcontents
\newpage

% ============================================================
\section{Introduction}
The Z-transform is an important tool for analyzing discrete-time signals
and systems because it converts a time-domain sequence into a complex
frequency-domain representation. In addition to the algebraic expression
$X(z)$, the region of convergence is essential because it determines the
set of values of $z$ for which the transform exists and helps distinguish
between right-sided and left-sided sequences.

Convolution is another fundamental DSP operation. Linear convolution
models the output of an LTI system, while circular convolution is closely
related to periodic sequences and discrete Fourier transform processing.
This laboratory combines these theoretical concepts with practical
experiments in Scilab, including audio and image data processing.

% ============================================================
\section{Objectives}
The objectives of this laboratory are:
\begin{itemize}[leftmargin=1.5em]
    \item to analyze discrete-time signals using the Z-transform;
    \item to determine the region of convergence for different signal types;
    \item to investigate Scilab functions for audio processing;
    \item to investigate image processing using the IPCV toolbox;
    \item to derive and implement linear convolution using folding-and-shifting and matrix methods;
    \item to derive and implement circular convolution using folding-and-shifting and matrix methods;
    \item to compute and visualize the energy of the input and output sequences.
\end{itemize}

% ============================================================
\section{Methodology}

\subsection{Software Environment}
The practical exercises were implemented in Scilab 2026.1.0. Built-in
sound functions were used for audio processing. For image processing,
the IPCV toolbox was installed and loaded through the ATOMS package
manager. The convolution algorithms were implemented directly using
loops, matrix operations, and discrete-time plots.

\subsection{Z-transform Definition}
For a discrete-time signal $x(n)$, the bilateral Z-transform is
\[
X(z)=\sum_{n=-\infty}^{\infty} x(n)z^{-n}.
\]
The ROC contains all values of $z$ for which this sum converges.
For a right-sided exponential sequence,
\[
\mathcal{Z}\{a^n u(n)\}=\frac{1}{1-a z^{-1}},
\qquad |z|>|a|.
\]
For a left-sided exponential sequence,
\[
\mathcal{Z}\{a^n u(-n-1)\}=-\frac{1}{1-a z^{-1}},
\qquad |z|<|a|.
\]

% ============================================================
\section{Experimental Content}

% ------------------------------------------------------------
\subsection{Exercise 1 -- Z-transform of a Finite-Duration Signal}
The signal is
\[
x(n)=2\delta(n+2)-\delta(n+1)+2\delta(n)-3\delta(n-1)+4\delta(n-2).
\]
Using the shifting property of the impulse,
\[
\mathcal{Z}\{\delta(n-k)\}=z^{-k},
\]
we obtain
\[
\boxed{X(z)=2z^2-z+2-3z^{-1}+4z^{-2}}.
\]

The sequence is finite-duration and contains nonzero samples for both
negative and positive values of $n$. Therefore, both $z=0$ and
$z=\infty$ are excluded from the ROC:
\[
\boxed{0<|z|<\infty}.
\]

% ------------------------------------------------------------
\subsection{Exercise 2 -- Sum of Two Right-Sided Exponentials}
The signal is
\[
x(n)=(0.5)^n u(n)+(0.4)^n u(n).
\]
The Z-transform is
\[
X(z)=\frac{1}{1-0.5z^{-1}}+\frac{1}{1-0.4z^{-1}}.
\]
Combining the two terms gives
\[
\boxed{
X(z)=\frac{2-0.9z^{-1}}
{1-0.9z^{-1}+0.2z^{-2}}
}.
\]
The individual ROCs are $|z|>0.5$ and $|z|>0.4$. Their intersection is
\[
\boxed{|z|>0.5}.
\]

% ------------------------------------------------------------
\subsection{Exercise 3 -- Two-Sided Exponential Signal}
The signal is
\[
x(n)=(0.5)^n u(n)+(0.9)^n u(-n-1).
\]
For the right-sided component,
\[
X_1(z)=\frac{1}{1-0.5z^{-1}},
\qquad |z|>0.5.
\]
For the left-sided component,
\[
X_2(z)=-\frac{1}{1-0.9z^{-1}},
\qquad |z|<0.9.
\]
Hence,
\[
\boxed{
X(z)=\frac{1}{1-0.5z^{-1}}
-\frac{1}{1-0.9z^{-1}}
}
\]
and the overall ROC is the intersection
\[
\boxed{0.5<|z|<0.9}.
\]
This annular ROC is characteristic of a two-sided sequence.

% ------------------------------------------------------------
\subsection{Exercise 4 -- Audio Processing in Scilab}
Scilab provides built-in functions for reading, playing, and analyzing
audio signals. In this experiment, the built-in \texttt{chimes.wav}
file was read using \texttt{wavread}. The first channel was selected for
analysis, the time-domain waveform was plotted, and the FFT was used to
obtain a magnitude spectrum.

The experiment reported:
\begin{itemize}
    \item sampling frequency $F_s=22050$ Hz;
    \item bit depth = 16 bits;
    \item number of channels = 2;
    \item samples per channel = 13921.
\end{itemize}

The waveform shows a short audio signal with several transient regions
whose amplitudes decay with time. The magnitude spectrum contains several
dominant peaks in the low-frequency range, which correspond to the main
frequency components of the audio signal.

\reportfigure[0.96\textwidth]{figures/ex4_audio.png}
{Audio waveform, magnitude spectrum, and audio-file information obtained in Scilab.}
{fig:ex4audio}

\subsubsection*{Scilab source}
\lstinputlisting[caption={Audio-processing script for Exercise 4.}]
{code/ex4_audio_demo.sce}

% ------------------------------------------------------------
\subsection{Exercise 5 -- Image Processing with IPCV}
The IPCV toolbox was used to demonstrate basic digital image processing.
After the toolbox was loaded, the sample image \texttt{baboon.png} was
read using \texttt{imread}. The original RGB image has size
$512\times512\times3$. It was converted to a grayscale image of size
$512\times512$ using \texttt{rgb2gray}. A histogram of grayscale
intensities was then generated.

The experiment demonstrates three common image-processing operations:
reading an image, converting between color representations, and
examining the intensity distribution with a histogram. The histogram
shows that the image uses a broad range of gray levels, with relatively
large populations in the middle and upper intensity ranges.

\reportfigure[0.96\textwidth]{figures/ex5_image.png}
{IPCV demonstration: original image, grayscale conversion, and grayscale histogram.}
{fig:ex5image}

\subsubsection*{Scilab source}
\lstinputlisting[caption={Image-processing script for Exercise 5.}]
{code/ex5_image_demo.sce}

% ------------------------------------------------------------
\subsection{Exercise 6 -- Linear Convolution}
Let
\[
x(n)=[1,2,-3,2,1], \qquad h(n)=[1,0,-1],
\]
with the first element of each sequence located at $n=0$.

\subsubsection{Convergence of the Convolution Sum}
The linear convolution is
\[
y(n)=\sum_{k=-\infty}^{\infty} h(k)x(n-k).
\]

If both $h(n)$ and $x(n)$ are causal, then $h(k)=0$ for $k<0$, and
$x(n-k)=0$ when $k>n$. Hence, for $n\ge0$,
\[
\boxed{y(n)=\sum_{k=0}^{n} h(k)x(n-k)},
\]
while $y(n)=0$ for $n<0$. For every fixed $n$, only finitely many terms
are nonzero.

If the system has a finite impulse response satisfying
\[
h(n)=0,\qquad n<0\ \text{or}\ n>M,
\]
then only $k=0,1,\ldots,M$ contribute to the convolution. Therefore,
\[
\boxed{y(n)=\sum_{k=0}^{M} h(k)x(n-k)}.
\]
Again, the convolution contains only $M+1$ terms for each $n$.

\subsubsection{Folding-and-Shifting Method}
For the given impulse response,
\[
h(0)=1,\qquad h(1)=0,\qquad h(2)=-1.
\]
Therefore,
\[
y(n)=x(n)-x(n-2).
\]
Evaluating the sum for $n=0,\ldots,6$ gives
\[
\boxed{y(n)=[1,2,-4,0,4,-2,-1]}.
\]

The signal energies are
\[
E_x=\sum_n |x(n)|^2
=1^2+2^2+(-3)^2+2^2+1^2
=\boxed{19},
\]
\[
E_y=\sum_n |y(n)|^2
=\boxed{42}.
\]

\reportfigure[0.96\textwidth]{figures/ex6_folding_shifting.png}
{Linear convolution using the folding-and-shifting implementation. The figure also shows the energy per sample.}
{fig:ex6fold}

\subsubsection*{Scilab source -- folding and shifting}
\lstinputlisting[caption={Linear convolution using folding and shifting.}]
{code/ex6_linear_folding_shifting.sce}

\subsubsection{Matrix Method}
A convolution matrix can be constructed from $x(n)$:
\[
\mathbf{X}=
\begin{bmatrix}
1&0&0\\
2&1&0\\
-3&2&1\\
2&-3&2\\
1&2&-3\\
0&1&2\\
0&0&1
\end{bmatrix},
\qquad
\mathbf{h}=
\begin{bmatrix}
1\\0\\-1
\end{bmatrix}.
\]
Then
\[
\mathbf{y}=\mathbf{X}\mathbf{h}
=
\begin{bmatrix}
1\\2\\-4\\0\\4\\-2\\-1
\end{bmatrix}.
\]
Thus the matrix method gives the same result as the folding-and-shifting
method:
\[
\boxed{y(n)=[1,2,-4,0,4,-2,-1]}.
\]

\reportfigure[0.72\textwidth]{figures/ex6_matrix.png}
{Linear-convolution matrix and output obtained with the matrix method.}
{fig:ex6matrix}

\subsubsection*{Scilab source -- matrix method}
\lstinputlisting[caption={Linear convolution using the matrix method.}]
{code/ex6_linear_matrix.sce}

% ------------------------------------------------------------
\subsection{Exercise 7 -- Circular Convolution}
For a periodic sequence $x(n)$ with period $N$,
\[
x(n)=x(n+qN),\qquad q\in\mathbb{Z}.
\]
Therefore, every integer index $n-k$ can be mapped into the interval
$0,1,\ldots,N-1$ using its remainder modulo $N$:
\[
x(n-k)=x[(n-k)\bmod N].
\]
If $h(k)$ is nonzero only for $k=0,1,\ldots,M-1$, the convolution over
one period is consequently
\[
\boxed{
y(n)=\sum_{k=0}^{M-1} h(k)x[(n-k)\bmod N]
}.
\]
This explains the circular-convolution expression used in the exercise.

For
\[
x(n)=[1,2,-3,2,1],
\]
\[
h(n)=[1,0,-1,-1,1],
\]
and $N=5$, both sequences have length five.

\subsubsection{Folding-and-Shifting Method}
Using the modulo index in the circular convolution sum gives
\[
\boxed{y(n)=[4,-4,-3,0,3]}.
\]
The energies are
\[
\boxed{E_x=19}, \qquad \boxed{E_y=50}.
\]

In Scilab, \texttt{pmodulo} is used rather than \texttt{modulo} so that
negative values of $n-k$ are mapped to the nonnegative range
$0,1,\ldots,N-1$. This avoids invalid array indices.

\reportfigure[0.96\textwidth]{figures/ex7_folding_shifting.png}
{Circular convolution using folding and shifting. The console confirms $y(n)=[4,-4,-3,0,3]$, $E_x=19$, and $E_y=50$.}
{fig:ex7fold}

\subsubsection*{Scilab source -- folding and shifting}
\lstinputlisting[caption={Circular convolution using folding and shifting.}]
{code/ex7_circular_folding_shifting.sce}

\subsubsection{Matrix Method}
The circular-convolution matrix generated from $h(n)$ is
\[
\mathbf{H}_c=
\begin{bmatrix}
1&1&-1&-1&0\\
0&1&1&-1&-1\\
-1&0&1&1&-1\\
-1&-1&0&1&1\\
1&-1&-1&0&1
\end{bmatrix}.
\]
With
\[
\mathbf{x}=
\begin{bmatrix}
1\\2\\-3\\2\\1
\end{bmatrix},
\]
we obtain
\[
\mathbf{y}=\mathbf{H}_c\mathbf{x}
=
\begin{bmatrix}
4\\-4\\-3\\0\\3
\end{bmatrix}.
\]
Thus, the matrix method verifies the result of the folding-and-shifting
method exactly.

\reportfigure[0.72\textwidth]{figures/ex7_matrix.png}
{Circular-convolution matrix and output obtained with the matrix method.}
{fig:ex7matrix}

\subsubsection*{Scilab source -- matrix method}
\lstinputlisting[caption={Circular convolution using the matrix method.}]
{code/ex7_circular_matrix.sce}

% ============================================================
\section{Discussion}
The theoretical Z-transform exercises show that the ROC is an essential
part of the transform description. Exercise 2 is purely right-sided and
therefore has an ROC outside its outermost pole, whereas Exercise 3 is
two-sided and has an annular ROC between the two pole magnitudes.

The audio and image experiments demonstrate that Scilab can process
real digital data in addition to numerical sequences. The audio example
combines time-domain inspection with frequency-domain analysis through
the FFT, while the image example uses IPCV to perform grayscale
conversion and histogram analysis.

For Exercise 6, both linear-convolution implementations produced
\([1,2,-4,0,4,-2,-1]\), confirming the correctness of the folding-and-
shifting and matrix approaches. For Exercise 7, both circular-
convolution implementations produced \([4,-4,-3,0,3]\). The use of
\texttt{pmodulo} is important in Scilab because it maps negative circular
indices into valid nonnegative remainders before conversion to Scilab's
one-based array indexing.

Overall, all required theoretical results and practical Scilab
experiments were completed successfully, and the numerical outputs from
independent methods agree with each other.

% ============================================================
\section{References}
\begin{enumerate}
    \item Digital Signal Processing, Lab \#4 handout, course material.
    \item Scilab Help System, sound-processing functions including \texttt{wavread}, \texttt{fft}, and \texttt{playsnd}.
    \item IPCV -- Image Processing and Computer Vision Toolbox for Scilab.
\end{enumerate}

\end{document}
